\documentclass[aspectratio=169,11pt]{beamer}

%% Shared preamble for Programming and Quantitative Methods in Finance Beamer presentations
%% Adapted from SPM-PEM/spmcourse-beamer.sty (English localisation)

\usepackage{xcolor}
\usepackage{amsmath,amssymb,bm,mathtools}
\usepackage{booktabs}
\usepackage{hyperref}
\usepackage[english]{babel}

\setbeamertemplate{footline}[frame number]
\setbeamertemplate{itemize items}[circle]
\setbeamercolor{structure}{fg=red!30!green!60!black}
\setbeamercolor{item}{fg=red!30!green!60!black}
\setbeamercolor{itemize item}{fg=red!30!green!60!black}

\definecolor{slidebg0}{RGB}{255,255,255}
\definecolor{slidebg1}{RGB}{220,220,220}
\definecolor{slidebg2}{RGB}{220,235,255}
\definecolor{slidebg3}{RGB}{220,255,230}
\definecolor{slidebg4}{RGB}{255,235,210}
\definecolor{slidebg5}{RGB}{235,220,255}
\definecolor{slidebg6}{RGB}{255,220,205}

\newcounter{slidebg}
\setcounter{slidebg}{1}
\newcommand{\alternateslidebg}{
  \ifnum\value{slidebg}=1
    \setbeamercolor{background canvas}{bg=slidebg1}\setcounter{slidebg}{2}
  \else\ifnum\value{slidebg}=2
    \setbeamercolor{background canvas}{bg=slidebg2}\setcounter{slidebg}{3}
  \else\ifnum\value{slidebg}=3
    \setbeamercolor{background canvas}{bg=slidebg3}\setcounter{slidebg}{4}
  \else\ifnum\value{slidebg}=4
    \setbeamercolor{background canvas}{bg=slidebg4}\setcounter{slidebg}{5}
  \else\ifnum\value{slidebg}=5
    \setbeamercolor{background canvas}{bg=slidebg5}\setcounter{slidebg}{6}
  \else
    \setbeamercolor{background canvas}{bg=slidebg6}\setcounter{slidebg}{1}
  \fi\fi\fi\fi\fi
}

\ExplSyntaxOn
\NewExpandableDocumentCommand{\tocsectionbg}{m}
  { slidebg\int_eval:n { \int_mod:nn { #1 } { 6 } + 1 } }
\ExplSyntaxOff

\setbeamertemplate{section in toc}{%
  \setbeamercolor{section in toc}{fg=black,bg=\tocsectionbg{\inserttocsectionnumber}}%
  \begin{beamercolorbox}[wd=\textwidth,sep=0.7ex,leftskip=1em,rightskip=1em]{section in toc}
    \inserttocsectionnumber.\enspace\inserttocsection
  \end{beamercolorbox}%
  \vspace{0.3ex}%
}
\newcommand{\newsection}[1]{\begin{frame}[plain]
\centering\vfill
{\usebeamercolor[fg]{structure}\LARGE\bfseries \insertsectionnumber\ #1}
\vfill\end{frame}}
\NewDocumentCommand{\bgsection}{s m}{%
  \alternateslidebg
  \section{#2}%
  \IfBooleanF{#1}{\newsection{#2}}%
}
\newcommand{\titleandoutline}{
  \alternateslidebg
  \frame{\titlepage}
  \begin{frame}{Outline}
  \tableofcontents[hideallsubsections]
  \end{frame}
}

\title{Programming and Quantitative Methods in Finance}
\subtitle{Králik Balázs (Office Hours Wed 15:30--17:00 E275) \\
Justin Huang (Office Hours Thu 11:30--13:00 E275)}
\author{}
\date{}

\begin{document}

\titleandoutline

\bgsection{Course structure}

\begin{frame}{Course structure}
  \begin{itemize}
    \item 24 classes in total.
    \item First half: 12 classes on Python programming (Králik Balázs).
    \item Second half: 12 classes on statistics (Justin Huang).
  \end{itemize}
\end{frame}

\begin{frame}{Assessment at a glance}
  \begin{itemize}
    \item Two midterms altogether: one in Python (week 6) and one in statistics
          (week 12). Each midterm is worth 30\%. Midterms: $2 \times 30\% = 60\%$
          of the final grade.
    \item 10 quizzes, five in each half. Each quiz is 4\% of the final grade.
          Quizzes: $10 \times 4\% = 40\%$ of the final grade.
    \item A quiz or midterm counts only when completed in class, closed book.
    \item There is one optional combined make-up exam in the exam period. If you
          take it, it overrides all previous scores for the final grade calculation,
          even if it is worse. However, it will not cause you to fail if your
          previous scores were sufficient to pass.
    \item The Python and statistics halves must each be passed separately in order
          to pass the course.
  \end{itemize}
\end{frame}

\bgsection{Getting started}

\begin{frame}{Installing Python}
  \begin{itemize}
    \item Install a current Python 3 release from \url{https://www.python.org/}
          or use a managed distribution such as Anaconda.
    \item During installation, make sure Python can be run from the
          command line. Note that on some systems, you may need to use
          \texttt{python3} or \texttt{Py} instead of \texttt{python}.
    \item Verify the installation with \texttt{python3 --version}.
    \item Use the same interpreter consistently for scripts, notebooks, and
          package installation.
  \end{itemize}
\end{frame}

\begin{frame}{Virtual environment}
  \begin{itemize}
    \item Keep course packages isolated from the system Python installation.
    \item From the repository root, create and activate \texttt{.venv}.
  \end{itemize}
  \begin{block}{macOS and Linux}
    \texttt{python3 -m venv .venv}\\
    \texttt{source .venv/bin/activate}\\
    \texttt{python -m pip install --upgrade pip}
  \end{block}
  \begin{block}{Windows PowerShell}
    \texttt{.venv\textbackslash Scripts\textbackslash Activate.ps1}
  \end{block}
\end{frame}

\begin{frame}{Integrated development environments (IDEs)}
  \begin{description}
    \item[Visual Studio Code] Recommended for this course: scripts, notebooks,
          extensions, and an integrated terminal. This is where exams will be given.
    \item[Jupyter] Excellent for exploratory analysis and presenting code,
          output, and plots together.
    \item[PyCharm] A full-featured Python IDE and a valid alternative.
    \item[Online editors] Useful for quick experiments, but less suitable for
          the complete course repository and local data files.
  \end{description}
\end{frame}

\begin{frame}{Visual Studio Code}
  \begin{enumerate}
    \item Install Visual Studio Code from \url{https://code.visualstudio.com/}.
    \item Install Microsoft's \textbf{Python} extension.
    \item Install Microsoft's \textbf{Jupyter} extension for notebook files.
    \item Open the course repository as a folder.
    \item Use ``Python: Select Interpreter'' and choose the course environment.
  \end{enumerate}
\end{frame}

\begin{frame}[fragile]{In-class task: Hello World}
  Now that Python and VS Code are installed, try your first program.
  \begin{block}{hello.py}
\begin{verbatim}
print("Hello World!")
\end{verbatim}
  \end{block}
  \begin{enumerate}
    \item Create a file called \texttt{hello.py} in the course folder.
    \item Run it with VS Code's ``Run Python File in Terminal'', or with
          \texttt{python3 hello.py} in a terminal.
    \item Confirm the output is exactly \texttt{Hello World!}
  \end{enumerate}
\end{frame}

\begin{frame}[fragile]{In-class task: comments and line continuation}
  \begin{enumerate}
    \item Add a comment above the \texttt{print} line explaining what the
          program does, and an inline comment after it.
  \end{enumerate}
  \begin{block}{Comments}
\begin{verbatim}
# Hello World program
print("Hello World!")  # this line displays the text
\end{verbatim}
  \end{block}
  \begin{enumerate}
    \setcounter{enumi}{1}
    \item Rewrite the \texttt{print} call using a backslash to split it over
          three lines, and run it again.
  \end{enumerate}
  \begin{block}{Line continuation}
\begin{verbatim}
print( \\
      "Hello World!" \\
     )
\end{verbatim}
  \end{block}
\end{frame}

\bgsection{Jupyter notebooks}

\begin{frame}{What is an \texttt{.ipynb} file?}
  \begin{itemize}
    \item A Jupyter notebook stores code cells, Markdown explanations, and
          output in one document.
    \item It is useful for data analysis, plots, and step-by-step demonstrations.
    \item It is mainly an educational and presentation format in this course;
          submitted classroom work may still be requested as a \texttt{.py} file.
    \item Notebook files are JSON documents internally, but should be edited
          through Jupyter or VS Code rather than by hand.
  \end{itemize}
\end{frame}

\begin{frame}{Running a notebook in VS Code}
  \begin{enumerate}
    \item Open the \texttt{.ipynb} file.
    \item Enable or install the Python and Jupyter extensions if prompted.
    \item Select the course \texttt{.venv} as the notebook kernel.
    \item Run cells individually, or run the notebook from the toolbar.
    \item Re-run relevant cells after changing code or input data.
  \end{enumerate}
  \begin{block}{Kernel reminder}
    If several Python installations exist, the first run may ask which one
    should execute the notebook. Choose the course environment.
  \end{block}
\end{frame}

\begin{frame}{Notebook troubleshooting}
  \begin{description}
    \item[No kernel] Select the \texttt{.venv} interpreter in the notebook's
          kernel picker.
    \item[Missing package] Install the package into the selected environment,
          then restart the kernel.
    \item[Stale output] Re-run the cell, or restart and run all cells in order.
    \item[Different results] Check the selected interpreter, package versions,
          working directory, and input data.
  \end{description}
\end{frame}

\begin{frame}{Notebook habits}
  \begin{itemize}
    \item Run cells in a meaningful order and keep the notebook reproducible.
    \item Explain non-obvious calculations in Markdown cells.
    \item Do not treat displayed output as current until its code has been run
          in the selected kernel.
    \item For submitted work, follow the file format and submission instructions
          given in class.
  \end{itemize}
\end{frame}

\bgsection{Python}

\begin{frame}{How the Python classes work}
  \begin{itemize}
    \item The first part introduces the week's concepts through explanations
          and examples.
    \item The second part is devoted to independent in-class problem solving.
    \item The emphasis is practical: students learn to find and fix syntax,
          implementation, and logical errors while working.
    \item Questions and productive collaboration are encouraged, provided that
          other students can work undisturbed.
  \end{itemize}
\end{frame}

\begin{frame}{Python half: 12 classes}
  \small
  \begin{tabular}{cccp{0.68\textwidth}}
    \toprule
    Week & Class & Quiz & Main topic \\
    \midrule
    1 & 1  &           & Basics, variables, and strings \\
      & 2  & $\checkmark$ & Numeric types, conversion, booleans, input, and conditionals \\
    2 & 3  &           & Sequence types, lists, and loops \\
      & 4  & $\checkmark$ & Functions \\
    3 & 5  &           & Review \\
      & 6  & $\checkmark$ & Modules, dates, and the \texttt{math} module \\
    4 & 7  &           & Pandas I \\
      & 8  & $\checkmark$ & Pandas II \\
    5 & 9  &           & Data visualization \\
      & 10 & $\checkmark$ & Monte Carlo simulations \\
    6 & 11 &          & Practice and preparation for the midterm \\
      & 12 &          & Midterm \\
    \bottomrule
  \end{tabular}
\end{frame}

\bgsection{Statistics}

\begin{frame}
  \small
  \begin{tabular}{cccp{0.80\textwidth}}
    \toprule
    Week & Class & Quiz & Main topic \\
    \midrule
    7 & 13 & & R essentials: VS Code, objects, vectors, data frames, CSV, missing data \\
      & 14 & $\checkmark$ & Descriptive statistics: mean, median, quartiles, IQR, variance, SD, plots \\
    8 & 15 & & Probability distributions: Normal, z scores, quantiles, random variables, simulation \\
      & 16 & $\checkmark$ & Sampling distributions, standard error, sample size, CLT \\
    9 & 17 & & Confidence intervals: estimates, SE, 95\% interpretation, margin of error \\
      & 18 & $\checkmark$ & Hypothesis testing I: p values, Type I error, t tests \\
    10 & 19 & & Hypothesis testing II: contingency tables, chi-square, one-way ANOVA \\
       & 20 & $\checkmark$ & Regression I: scatterplots, covariance, Pearson correlation, fitted line \\
    11 & 21 & & Regression II: residuals, $R^2$, SE, p values, prediction \\
       & 22 & $\checkmark$ & Inference reliability: CLT, bootstrap, uncertainty, power, multiple testing \\
    12 & 23 & & Integrated analysis: finance case, exploration, method choice, interpretation \\
       & 24 & & Statistics examination: 90 minutes, method selection, code/output interpretation \\
    \bottomrule
  \end{tabular}
\end{frame}

\begin{frame}{Statistics quizzes}
  \begin{itemize}
    \item 8--10 minutes at the start of the listed class.
    \item Previously taught material only.
    \item Quiz 1: R objects, data frames, variable types, missing values, basic data operations.
    \item Quiz 2: descriptive summaries, plots, z scores, Normal probabilities, simulation.
    \item Quiz 3: sampling variability, SD versus SE, confidence intervals, precision.
    \item Quiz 4: test selection, t tests, chi-square, ANOVA, p values.
    \item Quiz 5: Pearson $r$, regression, $R^2$, inference, residuals, prediction, causality.
    \item Optional extensions are not assessed.
  \end{itemize}
\end{frame}

\end{document}
